% Generated draft: requires mathematical and source review.
\documentclass[11pt]{article}
\usepackage[margin=25mm]{geometry}
\usepackage{fontspec}
\setmainfont{Latin Modern Roman}
\usepackage{amsmath,amssymb,mathtools,unicode-math}
\setmathfont{Latin Modern Math}
\usepackage{enumitem,needspace,xurl,hyperref}
\setlength{\parindent}{0pt}
\setlength{\parskip}{5pt}
\setlength{\emergencystretch}{3em}
\newcommand{\N}{\mathbb{N}}
\newcommand{\Z}{\mathbb{Z}}
\newcommand{\Q}{\mathbb{Q}}
\newcommand{\R}{\mathbb{R}}
\newcommand{\C}{\mathbb{C}}
\newcommand{\F}{\mathbb{F}}
\newcommand{\A}{\mathbb{A}}
\newcommand{\PP}{\mathbb{P}}
\newcommand{\E}{\mathbb{E}}
\newcommand{\eps}{\varepsilon}
\newcommand{\dd}{\,\mathrm{d}}
\begin{document}
\section*{UnsolvedMath ID: 700752}
\textbf{Source problem code:} catalogue entry
\subsection{Definitions and mathematical statement}
Let \(\Gamma=(V,E)\) be a finite simple undirected graph. Its valence is \(57\) if every vertex has exactly \(57\) neighbours. Here a Moore graph of valence \(57\) and diameter \(2\) means a connected graph with these parameters attaining the Moore bound
\[
|V|=1+57+57(57-1)=3250.
\]
An automorphism of \(\Gamma\) is a bijection \(\sigma:V\to V\) preserving adjacency in both directions. It has order \(2\) if \(\sigma\ne\operatorname{id}_V\) and \(\sigma^2=\operatorname{id}_V\).

The question proposed by A.~A.~Makhn\"ev in 1990 is whether every Moore graph \(\Gamma\) of valence \(57\), diameter \(2\), and order \(3250\) satisfies
\[
\forall\sigma\in\operatorname{Aut}(\Gamma),\qquad
\sigma^2=\operatorname{id}_V\ \Longrightarrow\ \sigma=\operatorname{id}_V.
\]
The formulation is conditional: it does not assert that such a graph exists.

For context, a graph has rank \(3\) if its automorphism group acts transitively on vertices and the stabilizer of a vertex has exactly three orbits on the vertex set. The supplied record attributes to M.~Aschbacher the result that the hypothetical Moore graph in question cannot have rank \(3\). This is supplied background attribution, not an independently verified result or an additional hypothesis.
\subsection{Short English statement}
Must every hypothetical Moore graph with \(3250\) vertices and valence \(57\) have no automorphism of order \(2\)? The question quantifies over all such graphs without asserting their existence. The supplied registry labels the problem open.
\subsection{Sources}
\begin{itemize}
\item Local UnsolvedMath registry, permanent ID \(700752\), exact source problem code \texttt{catalogue entry}, title \emph{Kourovka Notebook Problem 11.60}. Supplied upstream revision: \texttt{d371ee917ac9647fbb06cd6fc529deec6f3dda82}.
\item \emph{The Kourovka Notebook (Unsolved Problems in Group Theory)}, problem 11.60, proposed by A.~A.~Makhn\"ev, 11th issue (1990); supplied edition reference: 21st issue (2026), arXiv:\texttt{1401.0300v48}. The supplied collection description names E.~I.~Khukhro and V.~D.~Mazurov as editors.
\item M.~Aschbacher, \emph{J. Algebra} \textbf{19}, no.~4 (1971), 538--540. The supplied record cites this reference for the assertion that the hypothetical graph is not a graph of rank \(3\); no article title is supplied.
\end{itemize}
\end{document}
